\documentclass{book}
\usepackage[backend=biber,natbib=true,style=authoryear]{biblatex}
\addbibresource{/home/nguyen/1_NQBH/math/bib.bib}
\usepackage[utf8]{vietnam}
\usepackage[utf8]{inputenc}
\usepackage{float}
\usepackage{graphicx}
\usepackage[colorlinks=true,linkcolor=blue,urlcolor=red,citecolor=magenta]{hyperref}
\usepackage{amsmath,amssymb,amsthm}
\allowdisplaybreaks
\numberwithin{equation}{section}
\newtheorem{assumption}{Assumption}[section]
\newtheorem{lemma}{Lemma}[section]
\newtheorem{corollary}{Corollary}[section]
\newtheorem{definition}{Definition}[section]
\newtheorem{proposition}{Proposition}[section]
\newtheorem{theorem}{Theorem}[section]
\newtheorem{notation}{Notation}[section]
\newtheorem{remark}{Remark}[section]
\newtheorem{example}{Example}[section]
\newtheorem{ques}{Question}[section]
\newtheorem{problem}{Problem}[section]
\newtheorem{conjecture}{Conjecture}[section]
\usepackage[left=0.5in,right=0.5in,top=1.5cm,bottom=1.5cm]{geometry}
\usepackage{fancyhdr}
\pagestyle{fancy}
\fancyhf{}
\addtolength{\headheight}{0pt} % obsolete
\lhead{\scriptsize \chaptername~\thechapter}
\rhead{\itshape \scriptsize \nouppercase{\leftmark}} %\nouppercase !
\renewcommand{\chaptermark}[1]{\markboth{#1}{}}
\cfoot{\thepage}
\def\labelitemii{$\circ$}

\title{Legends}
\author{Collector: Nguyen Quan Ba Hong}
\date{\today}

\begin{document}
\maketitle
\setcounter{secnumdepth}{6}
\setcounter{tocdepth}{6}
\tableofcontents

%------------------------------------------------------------------------------%

\part{\href{https://en.wikipedia.org/wiki/Mathematician}{Mathematicians}}

\chapter{\href{https://en.wikipedia.org/wiki/Mathematician}{Wikipedia/Mathematician}}

\textbf{Mathematician.} \href{https://en.wikipedia.org/wiki/Euclid}{Euclid} (holding \href{https://en.wikipedia.org/wiki/Calipers}{calipers}), Greek mathematician, known as the ``Father of Geometry''

\noindent
\textbf{Occupation.}
\begin{itemize}
    \item \textbf{Occupation type.} \href{https://en.wikipedia.org/wiki/Academic}{Academic}
\end{itemize}
\textbf{Description.}
\begin{itemize}
    \item \textbf{Competencies.} \href{https://en.wikipedia.org/wiki/Mathematics}{Mathematics}, \href{https://en.wikipedia.org/wiki/Analytical_skill}{analytical skills} and \href{https://en.wikipedia.org/wiki/Critical_thinking}{critical thinking} skills.
    \item \textbf{Education required.} \href{https://en.wikipedia.org/wiki/Doctoral_degree}{Doctoral degree}, occasionally \href{https://en.wikipedia.org/wiki/Master%27s_degree}{master's degree}.
    \item \textbf{Fields of employment.}
    \begin{itemize}
        \item universities,
        \item private corporations,
        \item financial industry,
        \item government
    \end{itemize}
    \item \textbf{Related jobs.} \href{https://en.wikipedia.org/wiki/Statistician}{statistician}, \href{https://en.wikipedia.org/wiki/Actuary}{actuary}.
\end{itemize}
A \textit{mathematician} is someone who uses an extensive knowledge of \href{https://en.wikipedia.org/wiki/Mathematics}{mathematics} in their work, typically to solve \href{https://en.wikipedia.org/wiki/Mathematical_problem}{mathematical problems}.

Mathematicians are concerned with \href{https://en.wikipedia.org/wiki/Number}{numbers}, \href{https://en.wikipedia.org/wiki/Data}{data}, \href{https://en.wikipedia.org/wiki/Quantity}{quantity}, \href{https://en.wikipedia.org/wiki/Mathematical_structure}{structure}, \href{https://en.wikipedia.org/wiki/Space}{space}, \href{https://en.wikipedia.org/wiki/Mathematical_model}{models}, and \href{https://en.wikipedia.org/wiki/Mathematics#Change}{change}.

\section{History}
For broader coverage of this topic, see \href{https://en.wikipedia.org/wiki/History_of_mathematics}{History of mathematics}.

%
1 of the earliest known mathematicians was \href{https://en.wikipedia.org/wiki/Thales_of_Miletus}{Thales of Miletus} (c. 624--c.546 BC); he has been hailed as the 1st true mathematician and the 1st known individual to whom a mathematical discovery has been attributed.[Boyer (1991), \textit{A History of Mathematics}, p. 43]

He is credited with the 1st use of deductive reasoning applied to geometry, by deriving 4 corollaries to \href{https://en.wikipedia.org/wiki/Thales%27_Theorem}{Thales' Theorem}.
    
%
The number of known mathematicians grew when \href{https://en.wikipedia.org/wiki/Pythagoras_of_Samos}{Pythagoras of Samos} (c. 582--c. 507 BC) established the \href{https://en.wikipedia.org/wiki/Pythagoreans}{Pythagorean School}, whose doctrine it was that mathematics ruled the universe and whose motto was ``All is number''.[Boyer 1991, ``\textit{Ionia and the Pythagoreans}'', p. 49]

It was the Pythagoreans who coined the term ``mathematics'', and with whom the study of mathematics for its own sake begins.

%
The 1st woman mathematician recorded by history was \href{https://en.wikipedia.org/wiki/Hypatia}{Hypatia} of Alexandria (AD 350--415).

She succeeded her father as Librarian at the Great Library and wrote many works on applied mathematics.

Because of a political dispute, the Christian community in Alexandria punished her, presuming she was involved, by stripping her naked and scraping off her skin with clamshells (some say roofing tiles).[``\textit{Ecclesiastical History}, Bk VI: Chap. 15''. Archived from the original on 2014-08-14. Retrieved 2014-11-19.]

%
Science and mathematics in the Islamic world during the Middle Ages followed various models and modes of funding varied based primarily on scholars.

It was extensive patronage and strong intellectual policies implemented by specific rulers that allowed scientific knowledge to develop in many areas.

Funding for translation of scientific texts in other languages was ongoing throughout the reign of certain caliphs,[Abattouy, M., Renn, J. \& Weinig, P., 2001. \textit{Transmission as Transformation: The Translation Movements in the Medieval East and West in a Comparative Perspective}. Science in Context, 14(1-2), 1-12.] and it turned out that certain scholars became experts in the works they translated and in turn received further support for continuing to develop certain sciences.

As these sciences received wider attention from the elite, more scholars were invited and funded to study particular sciences.

An example of a translator and mathematician who benefited from this type of support was \href{https://en.wikipedia.org/wiki/Al-Khawarizmi}{al-Khawarizmi}.

A notable feature of many scholars working under Muslim rule in medieval times is that they were often polymaths.

Examples include the work on \href{https://en.wikipedia.org/wiki/Optics}{optics}, maths and \href{https://en.wikipedia.org/wiki/Astronomy}{astronomy} of \href{https://en.wikipedia.org/wiki/Ibn_al-Haytham}{Ibn al-Haytham}.

%
The \href{https://en.wikipedia.org/wiki/Renaissance}{Renaissance} brought an increased emphasis on mathematics and science to Europe.

During this period of transition from a mainly feudal and ecclesiastical culture to a predominantly secular one, many notable mathematicians had other occupations: \href{https://en.wikipedia.org/wiki/Luca_Pacioli}{Luca Pacioli} (founder of \href{https://en.wikipedia.org/wiki/Accounting}{accounting}); \href{https://en.wikipedia.org/wiki/Niccol%C3%B2_Fontana_Tartaglia}{Niccolò Fontana Tartaglia} (notable engineer and bookkeeper); \href{https://en.wikipedia.org/wiki/Gerolamo_Cardano}{Gerolamo Cardano} (earliest founder of probability and binomial expansion); \href{https://en.wikipedia.org/wiki/Robert_Recorde}{Robert Recorde} (physician) and \href{https://en.wikipedia.org/wiki/Fran%C3%A7ois_Vi%C3%A8te}{François Viète}s (lawyer). 

%
As time passed, many mathematicians gravitated towards universities.

An emphasis on free thinking and experimentation had begun in Britain's oldest universities beginning in the 17th century at Oxford with the scientists \href{https://en.wikipedia.org/wiki/Robert_Hooke}{Robert Hooke} and \href{https://en.wikipedia.org/wiki/Robert_Boyle}{Robert Boyle}, and at Cambridge where \href{https://en.wikipedia.org/wiki/Isaac_Newton}{Isaac Newton} was \href{https://en.wikipedia.org/wiki/Lucasian_Professor_of_Mathematics}{Lucasian Professor of Mathematics \& Physics}.

Moving into the 19th century, the objective of universities all across Europe evolved from teaching the ``\textit{regurgitation of knowledge}'' to ``\textit{encourag[ing] productive thinking}.''[Röhrs, ``\textit{The Classical Idea of the University},'' Tradition and Reform of the University under an International Perspective p. 20]

In 1810, Humboldt convinced the King of Prussia to build a university in Berlin based on \href{https://en.wikipedia.org/wiki/Friedrich_Schleiermacher}{Friedrich Schleiermacher}'s liberal ideas; the goal was to demonstrate the process of the discovery of knowledge and to teach students to ``\textit{take account of fundamental laws of science in all their thinking}.''

Thus, seminars and laboratories started to evolve.[Rüegg, ``Themes'', \textit{A History of the University in Europe, Vol. III}, p. 5--6]

%
British universities of this period adopted some approaches familiar to the Italian and German universities, but as they already enjoyed substantial freedoms and \href{https://en.wikipedia.org/wiki/Autonomy}{autonomy} the changes there had begun with the \href{https://en.wikipedia.org/wiki/Age_of_Enlightenment}{Age of Enlightenment}, the same influences that inspired Humboldt.

The Universities of \href{https://en.wikipedia.org/wiki/University_of_Oxford}{Oxford} and \href{https://en.wikipedia.org/wiki/University_of_Cambridge}{Cambridge} emphasized the importance of \href{https://en.wikipedia.org/wiki/Research}{research}, arguably more authentically implementing Humboldt's idea of a university than even German universities, which were subject to state authority.[Rüegg, ``Themes'', \textit{A History of the University in Europe, Vol. III}, p. 12]

Overall, science (including mathematics) became the focus of universities in the 19th and 20th centuries.

Students could conduct research in \href{https://en.wikipedia.org/wiki/Seminars}{seminars} or \href{https://en.wikipedia.org/wiki/Laboratories}{laboratories} and began to produce doctoral theses with more scientific content.[Rüegg, ``Themes'', \textit{A History of the University in Europe, Vol. III}, p. 13]

According to Humboldt, the mission of the \href{https://en.wikipedia.org/wiki/University_of_Berlin}{University of Berlin} was to pursue scientific knowledge.[Rüegg, ``Themes'', \textit{A History of the University in Europe, Vol. III}, p. 16]

\textit{The German university system fostered professional, bureaucratically regulated scientific research performed in well-equipped laboratories, instead of the kind of research done by private and individual scholars in Great Britain and France}.[Rüegg, ``Themes'', \textit{A History of the University in Europe, Vol. III}, p. 17--18]

In fact, Rüegg asserts that the \textit{German system is responsible for the development of the modern research university because it focused on the idea of ``freedom of scientific research, teaching and study.''}[Rüegg, ``Themes'', \textit{A History of the University in Europe, Vol. III}, p. 31]

\section{Required education}
Mathematicians usually cover a breadth of topics within mathematics in their \href{https://en.wikipedia.org/wiki/Undergraduate_education}{undergraduate education}, and then proceed to specialize in topics of their own choice at the \href{https://en.wikipedia.org/wiki/Graduate-level}{graduate level}.

In some universities, a \href{https://en.wikipedia.org/wiki/Qualifying_exam}{qualifying exam} serves to test both the breadth and depth of a student's understanding of mathematics; the students, who pass, are permitted to work on a \href{https://en.wikipedia.org/wiki/Doctoral_dissertation}{doctoral dissertation}.

\section{Activities}
\textsf{\href{https://en.wikipedia.org/wiki/Emmy_Noether}{Emmy Noether}, mathematical theorist and teacher.}

\subsection{Applied mathematics}
Main article: \href{https://en.wikipedia.org/wiki/Applied_mathematics}{Applied mathematics}.

%
Mathematicians involved with solving problems with applications in real life are called \href{https://en.wikipedia.org/wiki/Applied_mathematician}{applied mathematicians}.

Applied mathematicians are mathematical scientists who, with their specialized knowledge and \href{https://en.wikipedia.org/wiki/Professional}{professional} methodology, approach many of the imposing problems presented in related scientific fields.

With professional focus on a wide variety of problems, theoretical systems, and localized constructs, applied mathematicians work regularly in the study and formulation of \href{https://en.wikipedia.org/wiki/Mathematical_models}{mathematical models}.

Mathematicians and applied mathematicians are considered to be 2 of the STEM (science, technology, engineering, and mathematics) careers.

%
The discipline of \href{https://en.wikipedia.org/wiki/Applied_mathematics}{applied mathematics} concerns itself with mathematical methods that are typically used in science, engineering, business, and industry; thus, ``applied mathematics'' is a \href{https://en.wikipedia.org/wiki/Mathematical_science}{mathematical science} with specialized knowledge.

The term ``applied mathematics'' also describes the professional specialty in which mathematicians work on problems, often concrete but sometimes abstract.

As professionals focused on problem solving, \textit{applied mathematicians} look into the \textit{formulation, study}, and \textit{use of mathematical models} in \href{https://en.wikipedia.org/wiki/Science}{science}, \href{https://en.wikipedia.org/wiki/Engineering}{engineering}, \href{https://en.wikipedia.org/wiki/Business}{business}, and other areas of mathematical practice.

\subsection{Abstract mathematics}
Main article: \href{https://en.wikipedia.org/wiki/Pure_mathematics}{Pure mathematics}.

%
\href{https://en.wikipedia.org/wiki/Pure_mathematics}{Pure mathematics} is mathematics that studies entirely abstract \href{https://en.wikipedia.org/wiki/Concept}{concepts}.

From the 18th century onwards, this was a recognized category of mathematical activity, sometimes characterized as \textit{speculative mathematics},[See for example titles of works by Thomas Simpson from the mid-18th century: \textit{Essays on Several Curious and Useful Subjects in Speculative and Mixed Mathematics, Miscellaneous Tracts on Some Curious and Very Interesting Subjects in Mechanics, Physical Astronomy and Speculative Mathematics. Chisholm}, Hugh, ed. (1911). ``Simpson, Thomas''. Encyclop\ae dia Britannica. 25 (11th ed.). Cambridge University Press. p. 135.] and at variance with the trend towards meeting the needs of \href{https://en.wikipedia.org/wiki/Navigation}{navigation}, \href{https://en.wikipedia.org/wiki/Astronomy}{astronomy}, \href{https://en.wikipedia.org/wiki/Physics}{physics}, \href{https://en.wikipedia.org/wiki/Economics}{economics}, \href{https://en.wikipedia.org/wiki/Engineering}{engineering}, and other applications.

%
Another insightful view put forth is that \textit{pure mathematics is not necessarily \href{https://en.wikipedia.org/wiki/Applied_mathematics}{applied mathematics}}: it is possible to study abstract entities w.r.t. their intrinsic nature, and not be concerned with how they manifest in the real world.[Andy Magid, \textit{Letter from the Editor, in Notices of the AMS}, Nov 2005, American Mathematical Society, p. 1173. [1] Archived 2016-03-03 at the Wayback Machine]

\textit{Even though the pure and applied viewpoints are distinct philosophical positions, in practice there is much overlap in the activity of pure and applied mathematicians}.

%
\textit{To develop accurate models for describing the real world, many applied mathematicians draw on tools and techniques that are often considered to be ``pure'' mathematics}.

\textit{On the other hand, many pure mathematicians draw on natural and social phenomena as inspiration for their abstract research}.

\subsection{Mathematics teaching}
Many professional mathematicians also engage in the teaching of mathematics.

Duties may include:
\begin{itemize}
    \item teaching university mathematics courses;
    \item supervising undergraduate and graduate research; and
    \item serving on academic committees.
\end{itemize}

\subsection{Consulting}
Many careers in mathematics outside of universities involve consulting.

E.g., actuaries assemble and analyze data to estimate the probability and likely cost of the occurrence of an event such as death, sickness, injury, disability, or loss of property.

Actuaries also address financial questions, including those involving the level of pension contributions required to produce a certain retirement income and the way in which a company should invest resources to maximize its return on investments in light of potential risk.

Using their broad knowledge, actuaries help design and price insurance policies, pension plans, and other financial strategies in a manner which will help ensure that the plans are maintained on a sound financial basis.

%
As another example, mathematical finance will derive and extend the \href{https://en.wikipedia.org/wiki/Mathematical_model}{mathematical} or \href{https://en.wikipedia.org/wiki/Numerical_analysis}{numerical} models without necessarily establishing a link to \textit{financial theory}, taking observed market prices as input.

Mathematical consistency is required, not compatibility with \textit{economic theory}.

Thus, e.g., while a financial economist might study the structural reasons why a company may have a certain \href{https://en.wikipedia.org/wiki/Share_price}{share price}, a financial mathematician may take the share price as a given, and attempt to use \href{https://en.wikipedia.org/wiki/Stochastic_calculus}{stochastic calculus} to obtain the corresponding value of \href{https://en.wikipedia.org/wiki/Derivative_(finance)}{derivatives} of the \href{https://en.wikipedia.org/wiki/Stock}{stock} (see: \href{https://en.wikipedia.org/wiki/Valuation_of_options}{Valuation of options}; \href{https://en.wikipedia.org/wiki/Financial_modeling#Quantitative_finance}{Financial modeling}).

\section{Occupations}
\textsf{In 1938 in the United States, mathematicians were desired as teachers, calculating machine operators, mechanical engineers, accounting auditor bookkeepers, and actuary statisticians.}

According to the \href{https://en.wikipedia.org/wiki/Dictionary_of_Occupational_Titles}{Dictionary of Occupational Titles} occupations in mathematics include the following.[``020 OCCUPATIONS IN MATHEMATICS''. \textit{Dictionary Of Occupational Titles}. Archived from the original on 2012-11-02. Retrieved 2013-01-20.]
\begin{itemize}
    \item Mathematician
    \item Operations-Research Analyst
    \item Mathematical Statistician
    \item Mathematical Technician
    \item \href{https://en.wikipedia.org/wiki/Actuary}{Actuary}
    \item Applied Statistician
    \item Weight Analyst
\end{itemize}

\section{Quotations about mathematicians}
The following are quotations about mathematicians, or by mathematicians.
\begin{quotation}
    ``\textit{A mathematician is a device for turning coffee into theorems}.'' - Attributed to both \href{https://en.wikipedia.org/wiki/Alfr%C3%A9d_R%C3%A9nyi}{Alfréd Rényi}[15] and \href{https://en.wikipedia.org/wiki/Paul_Erd%C5%91s}{Paul Erd\"os}
\end{quotation}

\begin{quotation}
    ``\textit{Die Mathematiker sind eine Art Franzosen; redet man mit ihnen, so übersetzen sie es in ihre Sprache, und dann ist es alsobald ganz etwas anderes}.''
    
    (Mathematicians are [like] a sort of Frenchmen; if you talk to them, they translate it into their own language, and then it is immediately something quite different.) - \href{https://en.wikipedia.org/wiki/Johann_Wolfgang_von_Goethe}{Johann Wolfgang von Goethe}[16]
\end{quotation}

\begin{quotation}
    ``\textit{Each generation has its few great mathematicians$\ldots$ and [the others'] research harms no one}.'' - Alfred W. Adler ($\sim$1930), ``\textit{Mathematics and Creativity}''[17]
\end{quotation}

\begin{quotation}
    ``In short, I never yet encountered the mere mathematician who could be trusted out of equal roots, or one who did not clandestinely hold it as a point of his faith that x squared + px was absolutely and unconditionally equal to q. Say to one of these gentlemen, by way of experiment, if you please, that you believe occasions may occur where x squared + px is not altogether equal to q, and, having made him understand what you mean, get out of his reach as speedily as convenient, for, beyond doubt, he will endeavor to knock you down.'' - \href{https://en.wikipedia.org/wiki/Edgar_Allan_Poe}{Edgar Allan Poe}, \textit{The purloined letter}
\end{quotation}

\begin{quotation}
    ``A mathematician, like a painter or poet, is a maker of patterns. If his patterns are more permanent than theirs, it is because they are made with ideas.'' - \href{https://en.wikipedia.org/wiki/G._H._Hardy}{G. H. Hardy}, \textit{A Mathematician's Apology}
\end{quotation}

\begin{quotation}
    ``\textit{Some of you may have met mathematicians and wondered how they got that way}.'' - \href{https://en.wikipedia.org/wiki/Tom_Lehrer}{Tom Lehrer}
\end{quotation}

\begin{quotation}
    ``\textit{It is impossible to be a mathematician without being a poet in soul}.'' - \href{https://en.wikipedia.org/wiki/Sofia_Kovalevskaya}{Sofia Kovalevskaya}
\end{quotation}

\begin{quotation}
    ``\textit{There are 2 ways to do great mathematics. The first is to be smarter than everybody else. The second way is to be stupider than everybody else - but persistent.}'' - \href{https://en.wikipedia.org/wiki/Raoul_Bott}{Raoul Bott}
\end{quotation}

\begin{quotation}
    ``\textit{Mathematics is the queen of the sciences and arithmetic the queen of mathematics}.'' - \href{https://en.wikipedia.org/wiki/Carl_Friedrich_Gauss}{Carl Friedrich Gauss} [18]
\end{quotation}

\section{Prizes in mathematics}
There is no Nobel Prize in mathematics, though sometimes mathematicians have won the Nobel Prize in a different field, such as economics.

Prominent prizes in mathematics include the \href{https://en.wikipedia.org/wiki/Abel_Prize}{Abel Prize}, the \href{https://en.wikipedia.org/wiki/Chern_Medal}{Chern Medal}, the \href{https://en.wikipedia.org/wiki/Fields_Medal}{Fields Medal}, the \href{https://en.wikipedia.org/wiki/Gauss_Prize}{Gauss Prize}, the \href{https://en.wikipedia.org/wiki/Frederic_Esser_Nemmers_Prize}{Nemmers Prize}, the \href{https://en.wikipedia.org/wiki/Balzan_Prize}{Balzan Prize}, the \href{https://en.wikipedia.org/wiki/Crafoord_Prize}{Crafoord Prize}, the \href{https://en.wikipedia.org/wiki/Shaw_Prize}{Shaw Prize}, the \href{https://en.wikipedia.org/wiki/Steele_Prize}{Steele Prize}, the \href{https://en.wikipedia.org/wiki/Wolf_Prize}{Wolf Prize}, the \href{https://en.wikipedia.org/wiki/Schock_Prize}{Schock Prize}, and the \href{https://en.wikipedia.org/wiki/Nevanlinna_Prize}{Nevanlinna Prize}.

%
The \href{https://en.wikipedia.org/wiki/American_Mathematical_Society}{American Mathematical Society}, \href{https://en.wikipedia.org/wiki/Association_for_Women_in_Mathematics}{Association for Women in Mathematics}, and other mathematical societies offer several prizes aimed at increasing the representation of women and minorities in the future of mathematics.

\section{Mathematical autobiographies}
Several well known mathematicians have written autobiographies in part to explain to a general audience what it is about mathematics that has made them want to devote their lives to its study.

These provide some of the best glimpses into what it means to be a mathematician.

The following list contains some works that are not autobiographies, but rather essays on mathematics and mathematicians with strong autobiographical elements.
\begin{itemize}
    \item \textit{The Book of My Life} - \href{https://en.wikipedia.org/wiki/Girolamo_Cardano}{Girolamo Cardano}[19]
    \item \href{https://en.wikipedia.org/wiki/A_Mathematician's_Apology}{A Mathematician's Apology} - \href{https://en.wikipedia.org/wiki/G.H._Hardy}{G.H. Hardy}[20]
    \item \href{https://en.wikipedia.org/wiki/A_Mathematician's_Miscellany}{A Mathematician's Miscellany} (republished as Littlewood's miscellany) - \href{https://en.wikipedia.org/wiki/J._E._Littlewood}{J. E. Littlewood}[Littlewood, J. E. (1990) [Originally \textit{A Mathematician's Miscellany} published in 1953], Béla Bollobás (ed.), \textit{Littlewood's miscellany}, Cambridge University Press, ISBN 0-521-33702 X]
    \item \textit{I Am a Mathematician} - \href{https://en.wikipedia.org/wiki/Norbert_Wiener}{Norbert Wiener}[Wiener, Norbert (1956), \textit{I Am a Mathematician / The Later Life of a Prodigy}, The M.I.T. Press, ISBN 0-262-73007-3]
    \item \textit{I Want to be a Mathematician} - \href{https://en.wikipedia.org/wiki/Paul_R._Halmos}{Paul R. Halmos}
    \item \textit{Adventures of a Mathematician} - \href{https://en.wikipedia.org/wiki/Stanislaw_Ulam}{Stanislaw Ulam}[Ulam, S. M. (1976), \textit{Adventures of a Mathematician}, Charles Scribner's Sons, ISBN 0-684-14391-7]
    \item \textit{Enigmas of Chance} - \href{https://en.wikipedia.org/wiki/Mark_Kac}{Mark Kac}[Kac, Mark (1987), \textit{Enigmas of Chance/An Autobiography}, University of California Press, ISBN 0-520-05986-7]
    \item \textit{Random Curves} - \href{https://en.wikipedia.org/wiki/Neal_Koblitz}{Neal Koblitz}
    \item \href{https://en.wikipedia.org/wiki/Edward_Frenkel#Love_and_Math}{\textit{Love and Math}} - \href{https://en.wikipedia.org/wiki/Edward_Frenkel}{Edward Frenkel}
    \item \textit{Mathematics Without Apologies} - \href{https://en.wikipedia.org/wiki/Michael_Harris_(mathematician)}{Michael Harris}[Harris, Michael (2015), \textit{Mathematics without apologies/portrait of a problematic vocation}, Princeton University Press, ISBN 978-0-691-15423-7]
\end{itemize}

\section{See also}
\begin{itemize}
    \item \href{https://en.wikipedia.org/wiki/Lists_of_mathematicians}{Lists of mathematicians}
    \item \href{https://en.wikipedia.org/wiki/Human_computer}{Human computer}
    \item \href{https://en.wikipedia.org/wiki/Mathematical_joke}{Mathematical joke}
    \item \href{https://en.wikipedia.org/wiki/A_Mathematician's_Apology}{A Mathematician's Apology}
    \item \href{https://en.wikipedia.org/wiki/Men_of_Mathematics}{Men of Mathematics} (book)
    \item \href{https://en.wikipedia.org/wiki/Mental_calculator}{Mental calculator}\hfill$\square$
\end{itemize}

%------------------------------------------------------------------------------%

\chapter{\href{https://en.wikipedia.org/wiki/Laurent_Schwartz}{Wikipedia/Laurent Schwartz}}

\textbf{Laurent Schwartz.}
\begin{itemize}
    \item \textbf{Born.} Mar 5, 1915. \href{https://en.wikipedia.org/wiki/Paris}{Paris}, France.
    \item \textbf{Died.} Jul 4, 2002 (aged 87). Paris, France.
    \item \textbf{Nationality.} French
    \item \textbf{Alma mater.} \href{https://en.wikipedia.org/wiki/%C3%89cole_Normale_Sup%C3%A9rieure}{Ecole Normale Sup\'erieure}
    \item \textbf{Known for.}
    \begin{itemize}
        \item \href{https://en.wikipedia.org/wiki/Distribution_(mathematics)}{Theory of Distributions}
        \item \href{https://en.wikipedia.org/wiki/Schwartz_kernel_theorem}{Schwartz kernel theorem}
        \item \href{https://en.wikipedia.org/wiki/Schwartz_space}{Schwartz space}
        \item \href{https://en.wikipedia.org/wiki/Schwartz-Bruhat_function}{Schwartz-Bruhat function}
        \item \href{https://en.wikipedia.org/wiki/Radonifying_function}{Radonifying operator}
        \item \href{https://en.wikipedia.org/wiki/Cylinder_set_measure}{Cylinder set measure}
    \end{itemize}
    \item \textbf{Awards.} \href{https://en.wikipedia.org/wiki/Fields_Medal}{Fields Medal} (1950)
\end{itemize}
\textbf{Scientific career.}
\begin{itemize}
    \item \textbf{Fields.} Mathematics
    \item \textbf{Institutions.}
    \begin{itemize}
        \item \href{https://en.wikipedia.org/wiki/University_of_Strasbourg}{University of Strashbourg}
        \item \href{https://en.wikipedia.org/wiki/University_of_Nancy}{University of Nancy}
        \item \href{https://en.wikipedia.org/wiki/University_of_Grenoble}{University of Grenoble}
        \item \href{https://en.wikipedia.org/wiki/%C3%89cole_Polytechnique}{\'Ecole Polytechnique}
        \item \href{https://en.wikipedia.org/wiki/Universit%C3%A9_de_Paris_VII}{Universit\'e de Paris VII}
    \end{itemize}
    \item \textbf{Doctoral advisor.} \href{https://en.wikipedia.org/wiki/Georges_Valiron}{Georges Valiron}
    \item \textbf{Doctoral students.}
    \begin{itemize}
        \item \href{https://en.wikipedia.org/wiki/Maurice_Audin}{Maurice Audin}
        \item \href{https://en.wikipedia.org/wiki/Georges_Glaeser}{Georges Glaeser}
        \item \href{https://en.wikipedia.org/wiki/Alexander_Grothendieck}{Alexander Grothendieck}
        \item \href{https://en.wikipedia.org/wiki/Jacques-Louis_Lions}{Jacques-Louis Lions}
        \item \href{https://en.wikipedia.org/wiki/Bernard_Malgrange}{Bernard Malgrange}
        \item \href{https://en.wikipedia.org/wiki/Andr%C3%A9_Martineau}{Andr\'e Martineau}
        \item \href{https://en.wikipedia.org/wiki/Bernard_Maurey}{Bernard Maurey}
        \item \href{https://en.wikipedia.org/wiki/Leopoldo_Nachbin}{Leopoldo Nachbin}
        \item \href{https://en.wikipedia.org/wiki/Henri_Hogbe_Nlend}{Henri Hogbe Nlend}
        \item \href{https://en.wikipedia.org/wiki/Gilles_Pisier}{Gilles Pisier}
        \item \href{https://en.wikipedia.org/wiki/Fran%C3%A7ois_Treves}{Fran\c{c}ois Treves}
    \end{itemize}
    \item \textbf{Influenced.} \href{https://en.wikipedia.org/wiki/Per_Enflo}{Per Enflo}
\end{itemize}
\textit{Laurent-Moïse Schwartz} (Mar 5, 1915 - Jul 4, 2002) was a French mathematician.

He pioneered the \href{https://en.wikipedia.org/wiki/Theory}{theory} of \href{https://en.wikipedia.org/wiki/Distribution_(mathematics)}{distributions}, which gives a well-defined meaning to objects such as the \href{https://en.wikipedia.org/wiki/Dirac_delta_function}{Dirac delta function}.

He was awarded the \href{https://en.wikipedia.org/wiki/Fields_Medal}{Fields Medal} in 1950 for his work on the \href{https://en.wikipedia.org/wiki/Distribution_(mathematics)}{theory of distributions}.

For several years he taught at the \href{https://en.wikipedia.org/wiki/%C3%89cole_polytechnique}{École polytechnique}.
    
\section{Biography}

\subsection{Family}
Laurent Schwartz came from a Jewish family of \href{https://en.wikipedia.org/wiki/Alsace}{Alsatian} origin, with a strong scientific background: his father was a well-known \href{https://en.wikipedia.org/wiki/Surgeon}{surgeon}, his uncle \href{https://en.wikipedia.org/wiki/Robert_Debr%C3%A9}{Robert Debré} (who contributed to the creation of \href{https://en.wikipedia.org/wiki/UNICEF}{UNICEF}) was a famous \href{https://en.wikipedia.org/wiki/Pediatrics}{pediatrician}, and his great-uncle-in-law, \href{https://en.wikipedia.org/wiki/Jacques_Hadamard}{Jacques Hadamard}, was a famous mathematician.

%
During his training at \href{https://en.wikipedia.org/wiki/Lyc%C3%A9e_Louis-le-Grand}{Lycée Louis-le-Grand} to enter the \href{https://en.wikipedia.org/wiki/%C3%89cole_Normale_Sup%C3%A9rieure}{École Normale Supérieure}, he fell in love with \href{https://en.wikipedia.org/wiki/Marie-H%C3%A9l%C3%A8ne_Schwartz}{Marie-Hélène Lévy}, daughter of the probabilist \href{https://en.wikipedia.org/wiki/Paul_L%C3%A9vy_(mathematician)}{Paul Lévy} who was then teaching at the \href{https://en.wikipedia.org/wiki/%C3%89cole_polytechnique}{École polytechnique}.
    
Later they would have 2 children, Marc-André and Claudine.

Marie-Hélène was gifted in mathematics as well, as she contributed to the geometry of singular analytic spaces and taught at the \href{https://en.wikipedia.org/wiki/Universit%C3%A9_Lille_Nord_de_France}{University of Lille}.

%
Angelo Guerraggio describes ``Mathematics, politics and butterflies'' as his ``3 great loves''.[1]

\subsection{Education}
According to his teachers, Schwartz was an exceptional student.

He was particularly gifted in Latin, Greek and mathematics.

1 of his teachers told his parents: ``\textit{Beware, some will say your son has a gift for languages, but he is only interested in the scientific and mathematical aspect of languages: he should become a mathematician}.''

%
In 1934, he was admitted at the École Normale Supérieure, and in 1937 he obtained the \href{https://en.wikipedia.org/wiki/Agr%C3%A9gation}{agrégation} (with rank 2).

\subsection{World War II}
As a man of \href{https://en.wikipedia.org/wiki/Trotskyism}{Trotskyist} affinities and \href{https://en.wikipedia.org/wiki/Jew}{Jewish} descent, life was difficult for Schwartz during \href{https://en.wikipedia.org/wiki/World_War_II}{World War II}.

He had to hide and change his identity to avoid being \href{https://en.wikipedia.org/wiki/Deportation}{deported} after Nazi Germany overran France.

He worked for the \href{https://en.wikipedia.org/wiki/University_of_Strasbourg}{University of Strasbourg} (which had been relocated in \href{https://en.wikipedia.org/wiki/Clermont-Ferrand}{Clermont-Ferrand} because of the war) under the name of Laurent-Marie Sélimartin, while Marie-Hélène used the name Lengé instead of Lévy.

Unlike other mathematicians at Clermont-Ferrand such as \href{https://en.wikipedia.org/wiki/Jacques_Feldbau}{Feldbau}, the couple managed to escape the Nazis.

\subsection{Later career}
Schwartz taught mainly at \href{https://en.wikipedia.org/wiki/%C3%89cole_Polytechnique}{École Polytechnique}, from 1958 to 1980.
    
At the end of the war, he spent one year in \href{https://en.wikipedia.org/wiki/Grenoble}{Grenoble} (1944), then in 1945 joined the University of \href{https://en.wikipedia.org/wiki/Nancy,_France}{Nancy} on the advice of \href{https://en.wikipedia.org/wiki/Jean_Delsarte}{Jean Delsarte} and \href{https://en.wikipedia.org/wiki/Jean_Dieudonn%C3%A9}{Jean Dieudonné}, where he spent 7 years.
    
He was both an influential researcher and teacher, with students such as \href{https://en.wikipedia.org/wiki/Bernard_Malgrange}{Bernard Malgrange}, \href{https://en.wikipedia.org/wiki/Jacques-Louis_Lions}{Jacques-Louis Lions}, \href{https://en.wikipedia.org/wiki/Fran%C3%A7ois_Bruhat}{François Bruhat} and \href{https://en.wikipedia.org/wiki/Alexander_Grothendieck}{Alexander Grothendieck}.
    
He joined the science faculty of the \href{https://en.wikipedia.org/wiki/University_of_Paris}{University of Paris} in 1952.

In 1958 he became a teacher at the \href{https://en.wikipedia.org/wiki/%C3%89cole_polytechnique}{École polytechnique} after having at 1st refused this position.
    
From 1961 to 1963 the École polytechnique suspended his right to teach, because of his having signed the \href{https://en.wikipedia.org/wiki/Manifesto_of_the_121}{Manifesto of the 121} about the \href{https://en.wikipedia.org/wiki/Algerian_war}{Algerian war}, a gesture not appreciated by Polytechnique's military administration.

However, Schwartz had a lasting influence on mathematics at the École polytechnique, having reorganized both teaching and research there.

In 1965 he established the \href{https://en.wikipedia.org/wiki/Centre_de_math%C3%A9matiques_Laurent-Schwartz}{Centre de mathématiques Laurent-Schwartz} (CMLS) as its 1st director.

%
In 1973 he was elected corresponding member of the \href{https://en.wikipedia.org/wiki/French_Academy_of_Sciences}{French Academy of Sciences}, and was promoted to full membership in 1975.

\section{Mathematical legacy}
In 1950 at the \href{https://en.wikipedia.org/wiki/International_Congress_of_Mathematicians}{International Congress of Mathematicians}, Schwartz was a plenary speaker[Schwartz, Laurent (1950). ``\textit{Théorie des noyaux}'' (PDF). In: Proceedings of the International Congress of Mathematicians, Cambridge, Massachusetts, U.S.A., Aug 30--Sep 6, 1950. vol. 1. pp. 220--230.] and was awarded the \href{https://en.wikipedia.org/wiki/Fields_Medal}{Fields Medal} for his work on \href{https://en.wikipedia.org/wiki/Distribution_(mathematics)}{distributions}.

He was the 1st French mathematician to receive the Fields medal.

Because of his sympathy for \href{https://en.wikipedia.org/wiki/Trotskyism}{Trotskyism}, Schwartz encountered serious problems trying to enter the United States to receive the medal; however, he was ultimately successful.

%
The theory of distributions clarified the (then) mysteries of the \href{https://en.wikipedia.org/wiki/Dirac_delta_function}{Dirac delta function} and \href{https://en.wikipedia.org/wiki/Heaviside_step_function}{Heaviside step function}.

It helps to extend the theory of \href{https://en.wikipedia.org/wiki/Fourier_transform}{Fourier transforms} and is now of critical importance to the theory of \href{https://en.wikipedia.org/wiki/Partial_differential_equation}{partial differential equations}.

\section{Popular science}
Throughout his life, Schwartz actively worked to promote science and bring it closer to the general audience.

Schwartz said:
\begin{quotation}\it
    ``What are mathematics helpful for? Mathematics are helpful for physics.
    
    Physics helps us make fridges.
    
    Fridges are made to contain spiny lobsters, and spiny lobsters help mathematicians who eat them and have hence better abilities to do mathematics, which are helpful for physics, which helps us make fridges which$\ldots$''[3]
\end{quotation}

\section{Entomology}
\textsf{\href{https://en.wikipedia.org/wiki/Clanis_schwartzi}{Clanis schwartzi} Paratype \href{https://en.wikipedia.org/wiki/MHNT}{MHNT}.}

His mother, who was passionate about natural science, passed on her taste for \href{https://en.wikipedia.org/wiki/Entomology}{entomology} to Laurent.

His personal collection of 20,000 \href{https://en.wikipedia.org/wiki/Lepidoptera}{Lepidoptera} specimens, collected during his various travels was bequeathed to the \href{https://en.wikipedia.org/wiki/National_Museum_of_Natural_History_(France)}{Muséum national d'histoire naturelle}), the \href{https://en.wikipedia.org/wiki/Mus%C3%A9e_des_Confluences}{Science Museum of Lyon}, the \href{https://en.wikipedia.org/wiki/Mus%C3%A9um_de_Toulouse}{Museum of Toulouse} and the Museo de Historia Natural Alcide d'Orbigny in \href{https://en.wikipedia.org/wiki/Cochabamba}{Cochabamba} (Bolivia).
    
Several species discovered by Schwartz bear his name.

\section{Personal ideology}
Apart from his scientific work, Schwartz was a well-known outspoken \href{https://en.wikipedia.org/wiki/Intellectual}{intellectual}.

As a young socialist influenced by \href{https://en.wikipedia.org/wiki/Leon_Trotsky}{Leon Trotsky}, Schwartz opposed the \href{https://en.wikipedia.org/wiki/Totalitarianism}{}{totalitarianism} of the \href{https://en.wikipedia.org/wiki/Soviet_Union}{Soviet Union}, particularly under \href{https://en.wikipedia.org/wiki/Joseph_Stalin}{Joseph Stalin}.

Schwartz ultimately rejected \href{https://en.wikipedia.org/wiki/Trotskyism}{Trotskyism} for \href{https://en.wikipedia.org/wiki/Democratic_socialism}{democratic socialism}.

%
On his religious views, Schwartz called himself an atheist.[4]

\section{Books}

\subsection{Research articles}
\begin{itemize}
    \item \textit{Œuvres scientifiques. I}.
    
    With a general introduction to the works of Schwartz by Claude Viterbo and an appreciation of Schwartz by Bernard Malgrange.
    
    With 1 DVD.
    
    Documents Mathématiques (Paris), 9. Société Mathématique de France, Paris, 2011. x+523 pp. ISBN 978-2-85629-317-1
    \begin{quotation}
        the 1st half of his works in analysis and partial differential equations.
        
        After a preface by Claude Viterbo, which includes a few photos, one will find a note by Schwartz himself about his works, followed by a few original documents (letters, course notes), a presentation by Bernard Malgrange of the theory of distributions for which Schwartz received the Fields Medal in 1950, and a selection of articles covering the period 1944--1954.
    \end{quotation}
    \item \textit{Œuvres scientifiques. II}.
    
    With an appreciation of Schwartz by Alain Guichardet.
    
    With 1 DVD.
    
    Documents Mathématiques (Paris), 10.
    
    Société Mathématique de France, Paris, 2011. x+507 pp. ISBN 978-2-85629-318-8
    \begin{quotation}
        the 2nd half of his works in analysis and partial differential equations.
        
        After a note by Alain Guichardet on Schwartz and his seminars, one will find a selection of articles covering the period 1954--1966.
    \end{quotation}
    \item \textit{Œuvres scientifiques. III}.
    
    With appreciations of Schwartz by Gilles Godefroy and Michel Émery.
    
    With 1 DVD.
    
    Documents Mathématiques (Paris), 11. Société Mathématique de France, Paris, 2011. x+619 pp. ISBN 978-2-85629-319-5
    \begin{quotation}
        his works on Banach space theory (1968--1987), introduced by Gilles Godefroy, and on probability theory (1970--1996), presented by Michel Émery, as well as some articles of a historical nature (1955--1994).
    \end{quotation}
\end{itemize}

\subsection{Technical books}
\begin{itemize}
    \item \textit{Analyse hilbertienne}. Collection Méthodes. Hermann, Paris, 1979. ii+297 pp. ISBN 2-7056-5897-1
    \item \textit{Application of distributions to the theory of elementary particles in quantum mechanics}. Gordon and Breach, New York, NY, 1968. 144pp. ISBN 9780677300900
    \item \textit{Cours d'analyse. 1}. 2nd edition. Hermann, Paris, 1981. xxix+830 pp. ISBN 2-7056-5764-9
    \item \textit{Cours d'analyse. 2}. 2nd edition. Hermann, Paris, 1981. xxiii+475+21+75 pp. ISBN 2-7056-5765-7
    \item [5] \textit{Étude des sommes d'exponentielles. 2ième éd}. Publications de l'Institut de Mathématique de l'Université de Strasbourg, V. Actualités Sci. Ind., Hermann, Paris 1959 151 pp.
    \item \textit{Geometry and probability in Banach spaces}. Based on notes taken by Paul R. Chernoff. Lecture Notes in Mathematics, 852. Springer-Verlag, Berlin-New York, 1981. x+101 pp. ISBN 3-540-10691-X
    \item \textit{Lectures on complex analytic manifolds}. With notes by M. S. Narasimhan. Reprint of the 1955 edition. Tata Institute of Fundamental Research Lectures on Mathematics and Physics, 4. Published for the Tata Institute of Fundamental Research, Bombay; by Springer-Verlag, Berlin, 1986. iv+182 pp. ISBN 3-540-12877-8
    \item \textit{Mathematics for the physical sciences}. Hermann, Paris; Addison-Wesley Publishing Co., Reading, Mass.-London-Don Mills, Ont. 1966 358 pp.
    \item \textit{Radon measures on arbitrary topological spaces and cylindrical measures}. Tata Institute of Fundamental Research Studies in Mathematics, No. 6. Published for the Tata Institute of Fundamental Research, Bombay by Oxford University Press, London, 1973. xii+393 pp.
    \item \textit{Semimartingales and their stochastic calculus on manifolds}. Edited and with a preface by Ian Iscoe. Collection de la Chaire Aisenstadt. Presses de l'Université de Montréal, Montreal, QC, 1984. 187 pp. ISBN 2-7606-0660-0
    \item \textit{Semi-martingales sur des variétés, et martingales conformes sur des variétés analytiques complexes}. Lecture Notes in Mathematics, 780. Springer, Berlin, 1980. xv+132 pp. ISBN 3-540-09749-X
    \item Les tenseurs. \textit{Suivi de ``Torseurs sur un espace affine'' by Y. Bamberger and J.-P. Bourguignon}. 2nd edition. Hermann, Paris, 1981. i+203 pp. ISBN 2-7056-1376-5
    \item [6] \textit{Théorie des distributions}. Publications de l'Institut de Mathématique de l'Université de Strasbourg, No. IX-X. Nouvelle édition, entiérement corrigée, refondue et augmentée. Hermann, Paris 1966 xiii+420 pp.
\end{itemize}

\subsection{Seminar notes}
\begin{itemize}
    \item \textit{Séminaire Schwartz in Paris 1953 bis 1961}. Online edition: [1]
\end{itemize}

\subsection{Popular books}
\begin{itemize}
    \item \textit{Pour sauver l'université.} Editions du Seuil, 1983. 122 pp. ISBN 2020065878
    \item \textit{A mathematician grappling with his century}. Translated from the 1997 French original by Leila Schneps. Birkhäuser Verlag, Basel, 2001. viii+490 pp. ISBN 3-7643-6052-6
\end{itemize}

\section{See also}
\begin{itemize}
    \item \href{https://en.wikipedia.org/wiki/Schwartz_distribution}{Schwartz distribution}
    \item \href{https://en.wikipedia.org/wiki/Schwartz_kernel_theorem}{Schwartz kernel theorem}
    \item \href{https://en.wikipedia.org/wiki/Schwartz_space}{Schwartz space}
    \item \href{https://en.wikipedia.org/wiki/Schwartz-Bruhat_function}{Schwartz-Bruhat function}
    \item \href{https://en.wikipedia.org/wiki/Nicolas_Bourbaki}{Nicolas Bourbaki}\hfill$\square$
\end{itemize}

%------------------------------------------------------------------------------%

\chapter{\href{https://en.wikipedia.org/wiki/Jacques-Louis_Lions}{Wikipedia/Jacques-Louis Lions}}

\textbf{Jacques-Louis Lions.}
\begin{itemize}
    \item \textbf{Born.} May 3, 1928. \href{https://en.wikipedia.org/wiki/Grasse}{Grasse}, \href{https://en.wikipedia.org/wiki/Alpes-Maritimes}{Alpes-Maritimes}, \href{https://en.wikipedia.org/wiki/France}{France}.
    \item \textbf{Died.} May 17, 2001 (aged 73)
    \item \textbf{Nationality.} French.
    \item \textbf{Alma mater.} \href{https://en.wikipedia.org/wiki/University_of_Nancy}{University of Nancy}.
    \item \textbf{Known for.} PDEs.
    \item \textbf{Awards.} \href{https://en.wikipedia.org/wiki/Japan_Prize}{Japan Prize} (1991).
\end{itemize}
\textbf{Scientific career.}
\begin{itemize}
    \item \textbf{Fields.} Mathematics.
    \item \textbf{Institutions.}
    \begin{itemize}
        \item \href{https://en.wikipedia.org/wiki/%C3%89cole_Polytechnique}{\'Ecole Polytechnique}
        \item \href{https://en.wikipedia.org/wiki/Coll%C3%A8ge_de_France}{Coll\`ege de France}
    \end{itemize}
    \item \textbf{\href{https://en.wikipedia.org/wiki/Doctoral_advisor}{Doctoral advisor}.} \href{https://en.wikipedia.org/wiki/Laurent_Schwartz}{Laurent Schwartz}
    \item \textbf{Doctoral students.}
    \begin{itemize}
        \item \href{https://en.wikipedia.org/wiki/Alain_Bensoussan}{Alain Bensoussan}
        \item \href{https://en.wikipedia.org/wiki/Jean-Michel_Bismut}{Jean-Michel Bismut}
        \item \href{https://en.wikipedia.org/wiki/Ha%C3%AFm_Brezis}{Ha\"im Brezis}
        \item \href{https://en.wikipedia.org/wiki/Erol_Gelenbe}{Erol Gelenbe}
        \item \href{https://en.wikipedia.org/wiki/Roland_Glowinski}{Roland Glowinski}
        \item \href{https://en.wikipedia.org/wiki/Roger_Temam}{Roger Temam}
    \end{itemize}
\end{itemize}
\textit{Jacques-Louis Lions} ([1] 3 May 1928 - May 17, 2001) was a French mathematician who made contributions to the theory of \href{https://en.wikipedia.org/wiki/Partial_differential_equation}{partial differential equations} and to \href{https://en.wikipedia.org/wiki/Stochastic_processes}{stochastic control}, among other areas.

He received the \href{https://en.wikipedia.org/wiki/Society_for_Industrial_and_Applied_Mathematics}{SIAM}'s \href{https://en.wikipedia.org/wiki/John_von_Neumann_Lecture}{John von Neumann Lecture} prize in 1986 and numerous other distinctions.[2][3]

Lions is listed as an \href{https://en.wikipedia.org/wiki/ISI_highly_cited_researcher}{ISI highly cited researcher}.[4] 

\section{Biography}
After being part of the French Résistance in 1943 and 1944, J.-L. Lions entered the \href{https://en.wikipedia.org/wiki/%C3%89cole_Normale_Sup%C3%A9rieure}{École Normale Supérieure} in 1947.
    
He was a professor of mathematics at the Université of Nancy, the Faculty of Sciences of Paris, and the \href{https://en.wikipedia.org/wiki/%C3%89cole_polytechnique}{École polytechnique}.
    
%
In 1966 he sent an invitation to \href{https://en.wikipedia.org/wiki/Gury_Marchuk}{Gury Marchuk}, the soviet mathematician to visit Paris.

This was hand delivered by \href{https://en.wikipedia.org/wiki/General_De_Gaulle}{General De Gaulle} during his visit to \href{https://en.wikipedia.org/wiki/Akademgorodok}{Akademgorodok} in June of that year.[5]

%
He joined the prestigious \href{https://en.wikipedia.org/wiki/Coll%C3%A8ge_de_France}{Collège de France} as well as the French Academy of Sciences in 1973.
    
In 1979, he was appointed director of the Institut National de la Recherche en Informatique et Automatique (\href{https://en.wikipedia.org/wiki/INRIA}{INRIA}), where he taught and promoted the use of numerical simulations using finite elements integration.

Throughout his career, Lions insisted on the \textit{use of mathematics in industry}, with a particular involvement in the French space program, as well as in domains such as energy and the environment.

This eventually led him to be appointed director of the Centre National d'Etudes Spatiales (\href{https://en.wikipedia.org/wiki/CNES}{CNES}) from 1984 to 1992.

%
Lions was elected President of the \href{https://en.wikipedia.org/wiki/International_Mathematical_Union}{International Mathematical Union} in 1991 and also received the \href{https://en.wikipedia.org/wiki/Japan_Prize}{Japan Prize} and the \href{https://en.wikipedia.org/wiki/Harvey_Prize}{Harvey Prize} that same year.[3]

In 1992, the \href{https://en.wikipedia.org/wiki/University_of_Houston}{University of Houston} awarded him an honorary doctoral degree.

He was elected president of the \href{https://en.wikipedia.org/wiki/French_Academy_of_Sciences}{French Academy of Sciences} in 1996 and was also a Foreign Member of the \href{https://en.wikipedia.org/wiki/Royal_Society}{Royal Society} (ForMemRS)[6] and numerous other foreign academies.[2][3]

%
He has left a considerable body of work, among this more than 400 scientific articles, 20 volumes of mathematics that were translated into English and Russian, and major contributions to several collective works, including the 4000 pages of the monumental \textit{Mathematical analysis and numerical methods for science and technology} (in collaboration with Robert Dautray), as well as the \textit{Handbook of numerical analysis} in 7 volumes (with \href{https://en.wikipedia.org/wiki/Philippe_G._Ciarlet}{Philippe G. Ciarlet}).

%
His son \href{https://en.wikipedia.org/wiki/Pierre-Louis_Lions}{Pierre-Louis Lions} is also a well-known mathematician who was awarded a \href{https://en.wikipedia.org/wiki/Fields_Medal}{Fields Medal} in 1994.[7]

Both father and son have received honorary doctorates from \href{https://en.wikipedia.org/wiki/Heriot-Watt_University}{Heriot-Watt University} in 1986 and 1995 respectively.[8]

\section{Books}
\begin{itemize}
    \item with Enrico Magenes: \textit{Problèmes aux limites non homogènes et applications}. 3 vols., 1968, 1970
    \item \textit{Contrôle optimal de systèmes gouvernés par des équations aux dérivées partielles}. 1968
    \item with L. Cesari: \textit{Quelques méthodes de résolution des problèmes aux limites non linéaires}. 1969
    \item with Roger Dautray: \textit{Mathematical analysis and numerical methods for science and technology}. 9 vols., 1984/5
    \item with Philippe Ciarlet: \textit{Handbook of numerical analysis}. 7 vols.
    \item with Alain Bensoussan, Papanicolaou: \textit{Asymptotic analysis of periodic structures}. North Holland 1978
    \item \textit{Controlabilité exacte, perturbations et stabilisation de systèmes distribués}[9]
    \item with John E. Lagnese: \textit{Modelling Analysis and Control of Thin Plates}.\hfill$\square$
\end{itemize}

%------------------------------------------------------------------------------%

\chapter{\href{https://en.wikipedia.org/wiki/Roger_Temam}{Wikipedia/Roger Temam}}

\textbf{Roger Meyer Temam.}
\begin{itemize}
    \item \textbf{Born.} May 19, 1940 (age 80)
    \item \textbf{Nationality.} French
    \item \textbf{Alma mater.} \href{https://en.wikipedia.org/wiki/University_of_Paris}{University of Paris}
    \item \textbf{Known for.} \href{https://en.wikipedia.org/wiki/Navier-Stokes_equations}{Navier-Stokes equations}
\end{itemize}
\textbf{Scientific career.}
\begin{itemize}
    \item \textbf{Fields.} \href{https://en.wikipedia.org/wiki/Applied_mathematics}{Applied mathematics.}
    \item \textbf{Institutions.}
    \begin{itemize}
        \item \href{https://en.wikipedia.org/wiki/Paris-Sud_University}{Paris-Sud University} (\href{https://en.wikipedia.org/wiki/Orsay}{Orsay}).
        \item \href{https://en.wikipedia.org/wiki/Indiana_University_Bloomington}{Indiana University}
    \end{itemize}
    \item \textbf{Doctoral advisor.} \href{https://en.wikipedia.org/wiki/Jacques-Louis_Lions}{Jacques-Louis Lions}
    \item \textbf{Doctoral students.}
    \begin{itemize}
        \item \href{https://en.wikipedia.org/wiki/Etienne_Pardoux}{Etienne Pardoux}
        \item \href{https://en.wikipedia.org/wiki/Denis_Serre}{Denis Serre}
    \end{itemize}
\end{itemize}
\textit{Roger Meyer Temam} (born May 19, 1940) is a French applied mathematician working in \href{https://en.wikipedia.org/wiki/Numerical_analysis}{numerical analysis}, \href{https://en.wikipedia.org/wiki/Nonlinear_partial_differential_equations}{nonlinear partial differential equations} and \href{https://en.wikipedia.org/wiki/Fluid_mechanics}{fluid mechanics}.

He graduated from the \href{https://en.wikipedia.org/wiki/University_of_Paris}{University of Paris} - the \href{https://en.wikipedia.org/wiki/Sorbonne}{Sorbonne} in 1967, completing a doctorate (\textit{thèse d'Etat}) under the direction of \href{https://en.wikipedia.org/wiki/Jacques-Louis_Lions}{Jacques-Louis Lions}.

He has published over 400 articles, as well as 12 (authored or co-authored) books.

\section{Scientific work}
The 1st work of Temam in his thesis dealt with the \textit{fractional steps method}.

Thereafter, ``he has continually explored and developed new directions and techniques'':[2]
\begin{itemize}
    \item \href{https://en.wikipedia.org/wiki/Calculus_of_variations}{calculus of variations}, and the notion of \textit{duality} (book \#7), developing the mathematical framework for \textit{discontinuous} (in displacement) \textit{solutions}; a concept later used for his works on the \textit{mathematical theory of plasticity} (book \#5);
    \item mathematical formulation of the equilibrium of a plasma in a cavity, expressed as a nonlinear \href{https://en.wikipedia.org/wiki/Free_boundary_problem}{free boundary problem};[R. Temam, A nonlinear eigenvalue problem: the shape at equilibrium of a confined plasma, \textit{Arch. Rational Mech. Anal.}, 60, 1975, 51--73.]
    \item \href{https://en.wikipedia.org/wiki/Korteweg-de_Vries_equation}{Korteweg-de Vries equation};[R. Temam, Sur un problème non linéaire, \textit{J. Math. Pures Appl.}, 48, 1969, 159--172.]
    \item \href{https://en.wikipedia.org/wiki/Kuramoto%E2%80%93Sivashinsky_equation}{Kuramoto-Sivashinsky equation};[5]
    \item \href{https://en.wikipedia.org/wiki/Euler_equations}{Euler equations} in a bounded domain;[R. Temam, On the Euler equations of incompressible perfect fluids, \textit{J. Funct. Anal.}, 20, 1975, 32--43.]
    \item infinite-dimensional \href{https://en.wikipedia.org/wiki/Dynamical_systems}{dynamical systems} theory.
    
    In particular, he studied the existence of the finite-dimensional global \href{https://en.wikipedia.org/wiki/Attractor}{attractor} for many dissipative equations of mathematical physics, including the incompressible \href{https://en.wikipedia.org/wiki/Navier-Stokes_equations}{Navier-Stokes equations}.[P. Constantin, C. Foias, O. Manley and R. Temam, Determining modes and fractal dimension of turbulent flows, \textit{J. Fluid Mech.}, 150, 1985, 427--440.][C. Foias, O.P. Manley and R. Temam, Physical estimates of the number of degrees of freedom in free convection, \textit{Phys. Fluids}, 29, 1986, 3101--3103.]
    
    He was also the co-founder of the notion of \textit{inertial manifolds}[C. Foias, G.R. Sell and R. Temam, Inertial manifolds for nonlinear evolutionary equations, \textit{J. Diff. Equ.}, 73, 1988, 309--353.] together with Ciprian Foias and \href{https://en.wikipedia.org/wiki/George_R._Sell}{George R. Sell} and of exponential attractors[A. Eden, C. Foias, B. Nicolaenko and R. Temam, \textit{Exponential attractors for dissipative evolution equations}, Collection Recherches en Mathématiques Appliquées, Masson, Paris, and John Wiley, England, 1994.] together with Alp Eden, Ciprian Foias and Basil Nicolaenko;[2]
    \item \href{https://en.wikipedia.org/wiki/Optimal_control}{optimal control} of the incompressible Navier-Stokes equations as a tool for the \textit{control of turbulence};[F. Abergel and R. Temam, On some control problems in fluid mechanics, \textit{Theoret. Comput. Fluid Dynamics}, 1, 1990, 303--325.]
    \item \href{https://en.wikipedia.org/wiki/Boundary_layer}{boundary layer} phenomena for incompressible flows.[12]
\end{itemize}
Temam's main activities concern the study of \href{https://en.wikipedia.org/wiki/Geophysical_fluid_dynamics}{geophysical flows}, the atmosphere and oceans.[2]

This started in the 1990s by collaboration with Jacques-Louis Lions and Shouhong Wang.[J.L. Lions, R. Temam and S. Wang, New formulations of the primitive equations of the atmosphere and applications, \textit{Nonlinearity}, 5, 1992, 237--288.][J.L. Lions, R. Temam and S. Wang, On the equations of the large-scale ocean, \textit{Nonlinearity}, 5, 1992, 1007--1053.][M. Coti Zelati, M. Frémond, R. Temam and J. Tribbia, Uniqueness, regularity and maximum principles for the equations of the atmosphere with humidity and saturation, \textit{Physica D}, 264, 2013, 49-65, https://doi.org/10.1016/j.physd.2013.08.007][Y. Cao, M. Hamouda, R. Temam, J. Tribbia and X. Wang, The equations of the multi-phase humid atmosphere expressed as a quasi variational inequality, \textit{Nonlinearity}, 31, 2018, 4692-4723, https://doi.org/10.1088/1361-6544/aad525.]

%
According to the \href{https://en.wikipedia.org/wiki/Mathematical_Genealogy_Project}{Mathematical Genealogy Project} database,[17][18] he holds the first position in the top 50 advisors.

More than 30 of his students are now full professors all over the world, and have themselves many descendants.[19]

\section{Administrative activities}
Temam became a professor at the \href{https://en.wikipedia.org/wiki/Paris-Sud_University}{Paris-Sud University} at Orsay in 1968.

There, he co-founded the Laboratory of Numerical and Functional Analysis which he directed from 1972 to 1988.

He was also a \textit{Maître de Conférences} at the \href{https://en.wikipedia.org/wiki/Ecole_Polytechnique}{Ecole Polytechnique} in Paris from 1968 to 1986.[20]

%
In 1983, Temam co-founded the French \href{https://en.wikipedia.org/wiki/Soci%C3%A9t%C3%A9_de_Math%C3%A9matiques_Appliqu%C3%A9es_et_Industrielles}{Société de Mathématiques Appliquées et Industrielles} (SMAI), analogous to the \href{https://en.wikipedia.org/wiki/Society_for_Industrial_and_Applied_Mathematics}{Society for Industrial and Applied Mathematics} (SIAM), and served as its 1st president.[21]
    
He was also 1 of the founders of the \href{https://en.wikipedia.org/wiki/International_Congress_on_Industrial_and_Applied_Mathematics}{International Congress on Industrial and Applied Mathematics} (ICIAM) series and was the chair of the steering committee of the 1st ICIAM meeting held in Paris in 1987; and the chair of the standing committee of the 2nd ICIAM meeting held in Washington, D.C., in 1991.[22]

He was the Editor-in-Chief of the mathematical journal M2AN[23] from 1986 to 1997.

%
Temam has been the Director of the Institute for Scientific Computing \& Applied Mathematics (ISCAM)[24] at \href{https://en.wikipedia.org/wiki/Indiana_University_Bloomington}{Indiana University} since 1986 (co-director with Ciprian Foias from 1986 to 1992).

He is also a College Professor (part-time till 2003) and he has been a Distinguished Professor since 2014.[25]

\section{Books}
\begin{enumerate}
    \item (with G.-M. Gie, M. Hamouda and C.-Y. Jung): \textit{Singular perturbations and boundary layers}, Springer-Verlag, New-York, 2018.
    \item (with A. Miranville): \textit{Mathematical Modelling in Continuum Mechanics}, Cambridge University Press, 2001. French Translation, Springer-Verlag France, 2002. Chinese Translation, Tsinghua University Press, 2004. 2nd English Edition 2005. Russian translation, Moskva Linom, 2013.
    \item (with T. Dubois and F. Jauberteau): \textit{Dynamic, multilevel methods and the numerical simulation of turbulence}; Cambridge University Press, 1999.
    \item \textit{Infinite Dimensional Dynamical Systems in Mechanics and Physics}, Springer-Verlag, New-York, Applied Mathematical Sciences Series, vol. 68, 1988. 2nd augmented edition, 1997. Reprinted in China by Beijing World Publishing Corp., 2000.
    \item \textit{Mathematical Problems in Plasticity}, Gauthier-Villars, Paris, 1983 (in French). English Transl., Gauthier-Villars, New-York, 1985. Russian Transl., Nauk, Moscow, 1991. ``Republished by Dover books in Physics, 2018.''
    \item \textit{Navier-Stokes Equations}, North-Holland Pub. Company, in English, 1977, 500 pages. Revised editions 1979, 1984 and 1985. Russian Translation, Mir, Moscow, 1981. ``Republished in the AMS-Chelsea Series, AMS, Providence, 2001.''
    \item (with I. Ekeland): \textit{Convex Analysis and Variational Problems}. Dunod, Paris, 1974, 350 pages (in French). English Translation, North-Holland, Amsterdam, 1976. Russian Translation, Mir, Moscow, 1979. ``English version republished in the Series 'Classics in Applied Mathematics', SIAM, Philadelphia, 1999.''
\end{enumerate}

\section{Awards \& honors}
\begin{itemize}
    \item Fellow of the \href{https://en.wikipedia.org/wiki/American_Academy_of_Arts_and_Sciences}{American Academy of Arts and Sciences} (2015),[26] of the \href{https://en.wikipedia.org/wiki/American_Mathematical_Society}{American Mathematical Society} (2013),[27] of the American Association for the Advancement of Science (2011),[28] of the Society for Industrial and Applied Mathematics (2009).[29]
    \item Knight of the \href{https://en.wikipedia.org/wiki/Legion_of_Honour}{Legion of Honor}, France, 2012.[30]
    \item Member of the \href{https://en.wikipedia.org/wiki/French_Academy_of_Sciences}{French Academy of Sciences} since 2007.[31]\hfill$\square$
\end{itemize}

%------------------------------------------------------------------------------%

\chapter{\href{https://en.wikipedia.org/wiki/Haim_Brezis}{Wikipedia/Ha\"im Brezis}}

\textbf{Haïm Brezis.}
\begin{itemize}
    \item \textbf{Born.} Jun 1m 1944 (age 76). Riom-ès-Montagnes, Cantal, France.
    \item \textbf{Nationality.} French
    \item \textbf{Alma mater.} \href{https://en.wikipedia.org/wiki/University_of_Paris}{University of Paris}
    \item \textbf{Known for.}
    \begin{itemize}
        \item \href{https://en.wikipedia.org/wiki/Brezis-Gallouet_inequality}{Brezis-Gallouet inequality}
        \item \href{https://en.wikipedia.org/wiki/Bony-Brezis_theorem}{Bony-Brezis theorem}
        \item \href{https://en.wikipedia.org/wiki/Brezis-Lieb_lemma}{Brezis-Lieb lemma}
    \end{itemize}
\end{itemize}
\textbf{Scientific career.}
\begin{itemize}
    \item \textbf{Fields.} Mathematics
    \item \textbf{Institutions.} \href{https://en.wikipedia.org/wiki/Pierre_and_Marie_Curie_University}{Pierre and Marie Curie University}
    \item \textbf{Doctoral advisor.}
    \begin{itemize}
        \item \href{https://en.wikipedia.org/wiki/Gustave_Choquet}{Gustave Choquet}
        \item \href{https://en.wikipedia.org/wiki/Jacques-Louis_Lions}{Jacques-Louis Lions}
    \end{itemize}
    \item \textbf{Doctoral students.}
    \begin{itemize}
        \item \href{https://en.wikipedia.org/wiki/Abbas_Bahri}{Abbas Bahri}
        \item \href{https://en.wikipedia.org/wiki/Jean-Michel_Coron}{Henri Berestycki}
        \item \href{https://en.wikipedia.org/wiki/Jean-Michel_Coron}{Jean-Michel Coron}
        \item \href{https://en.wikipedia.org/wiki/Jes%C3%BAs_Ildefonso_D%C3%ADaz}{Jes\'us Ildefonso D\'iaz}
        \item \href{https://en.wikipedia.org/wiki/Pierre-Louis_Lions}{Pierre-Louis Lions}
        \item \href{https://en.wikipedia.org/wiki/Juan_Luis_V%C3%A1zquez_Su%C3%A1rez}{Juan Luis V\'azquez Su\'arez}
    \end{itemize}
\end{itemize}
\textit{Haïm Brezis} (born Jun 1, 1944) is a French mathematician who works in \href{https://en.wikipedia.org/wiki/Functional_analysis}{functional analysis} and \href{https://en.wikipedia.org/wiki/Partial_differential_equation}{partial differential equations}.

\section{Biography}
Born in \href{https://en.wikipedia.org/wiki/Riom-%C3%A8s-Montagnes}{Riom-ès-Montagnes}, \href{https://en.wikipedia.org/wiki/Cantal}{Cantal}, France.
    
Brezis is the son of a Romanian immigrant father, who came to France in the 1930s, and a Jewish mother who fled from the Netherlands.

His wife, \href{https://en.wikipedia.org/wiki/Michal_Govrin}{Michal Govrin}, a native Israeli, works as a novelist, poet, and theater director.[1]

Brezis received his Ph.D. from the \href{https://en.wikipedia.org/wiki/University_of_Paris}{University of Paris} in 1972 under the supervision of \href{https://en.wikipedia.org/wiki/Gustave_Choquet}{Gustave Choquet}.

He is currently a Professor at the \href{https://en.wikipedia.org/wiki/Pierre_and_Marie_Curie_University}{Pierre and Marie Curie University} and a Visiting Distinguished Professor at \href{https://en.wikipedia.org/wiki/Rutgers_University}{Rutgers University}.

He is a member of the \href{https://en.wikipedia.org/wiki/Academia_Europaea}{Academia Europaea} (1988) and a foreign associate of the \href{https://en.wikipedia.org/wiki/United_States_National_Academy_of_Sciences}{United States National Academy of Sciences} (2003).

In 2012 he became a fellow of the \href{https://en.wikipedia.org/wiki/American_Mathematical_Society}{American Mathematical Society}.[2]

He holds honorary doctorates from several universities including \href{https://en.wikipedia.org/wiki/National_Technical_University_of_Athens}{National Technical University of Athens}.[3]

Brezis is listed as an \href{https://en.wikipedia.org/wiki/ISI_highly_cited_researcher}{ISI highly cited researcher}.[4]

\section{Works}
\begin{itemize}
    \item \textit{Opérateurs maximaux monotones et semi-groupes de contractions dans les espaces de Hilbert} (1973)
    \item \textit{Analyse Fonctionnelle}. Théorie et Applications (1983)
    \item \textit{Haïm Brezis. Un mathématicien juif}. Entretien Avec Jacques Vauthier. Collection Scientifiques \& Croyants. Editions Beauchesne, 1999. ISBN 978-2-7010-1335-0, ISBN 2-7010-1335-6
    \item \textit{Functional Analysis, Sobolev Spaces and Partial Differential Equations}, Springer; 1st Edition. edition (November 10, 2010), ISBN 978-0-387-70913-0, ISBN 0-387-70913-4
\end{itemize}

\section{See also}
\begin{itemize}
    \item \href{https://en.wikipedia.org/wiki/Bony%E2%80%93Brezis_theorem}{Bony-Brezis theorem}
    \item \href{https://en.wikipedia.org/wiki/Brezis%E2%80%93Gallouet_inequality}{Brezis-Gallouet inequality}
    \item \href{https://en.wikipedia.org/wiki/Brezis%E2%80%93Lieb_lemma}{Brezis-Lieb lemma}\hfill$\square$
\end{itemize}

%------------------------------------------------------------------------------%

\chapter{\href{https://en.wikipedia.org/wiki/Henri_Berestycki}{Wikipedia/Henri Berestycki}}

\textit{Henri Berestycki} (born Mar 25, 1951, in Paris)[1] is a French mathematician who obtained his PhD from Université Paris VI - \href{https://en.wikipedia.org/wiki/Universit\%C3\%A9_Pierre_et_Marie_Curie}{Université Pierre et Marie Curie} in 1975.

His Dissertation was titled \textit{Contributions à l'étude des problèmes elliptiques non linéaires}, and his doctoral advisor was \href{https://en.wikipedia.org/wiki/Haim_Brezis}{Haim Brezis}.[2]

He was an \href{https://en.wikipedia.org/wiki/Leonard_Eugene_Dickson}{L.E. Dickson} Instructor in Mathematics at the \href{https://en.wikipedia.org/wiki/University_of_Chicago}{University of Chicago} from 1975--77, after which he returned to France to continue his research.

He has made many contributions in \textit{nonlinear analysis}, ranging from \textit{nonlinear elliptic equations, hamiltonian systems, spectral theory of elliptic operators}, and with applications to the description of \textit{mathematical modelling of fluid mechanics and combustion}.

His current research interests include the mathematical modelling of financial markets, mathematical models in biology and especially in ecology, and modelling in social sciences (in particular, urban planning and criminology).

For these latter topics, he obtained an \href{http://erc.europa.eu/advanced-grants}{ERC Advanced grant} in 2012.

%
He worked at the French National Center of Scientific Research (\href{https://en.wikipedia.org/wiki/CNRS}{CNRS}), then moved to an appointment as Professor at Univ. Paris XIII (1983--1985).

He became a Professor of Mathematics in 1988 at Université Pierre et Marie Curie, Paris VI (1988--2001 of ``exceptional class'' since 1993), and became Professor at \href{https://en.wikipedia.org/wiki/Ecole_normale_sup%C3%A9rieure}{Ecole normale supérieure}, Paris (1994--1999), and part-time professor \href{https://en.wikipedia.org/wiki/Ecole_Polytechnique}{Ecole Polytechnique} (1987--1999).

He is also a visiting Professor in the Department of Mathematics at the University of Chicago, and was also co-director of the Stevanovich Center of Financial Mathematics in Chicago.

He is currently the Directeur d'études (Research Professor) at \href{https://en.wikipedia.org/wiki/School_for_Advanced_Studies_in_the_Social_Sciences}{École des hautes études en sciences sociales} (\href{https://en.wikipedia.org/wiki/EHESS}{EHESS}), since 2001.

\section{Services}
\begin{itemize}
    \item National Committee of French universities (1992--1995).
    \item Since 2002 director of Centre d'analyse et mathématique sociales (CAMS), CNRS -EHESS.
    \item Vice-president, EHESS (2004--2006).
    \item Member of the thesis prize committee of the universities of Paris (since 2006).
\end{itemize}

\section{Awards}
\begin{itemize}
    \item Carrière Prize(1988)
    \item \href{https://en.wikipedia.org/wiki/Sophie_Germain#Sophie_Germain_Prize}{Prix Sophie Germain} of the \href{https://en.wikipedia.org/wiki/French_Academy_of_Sciences}{French Academy of Sciences} (2004),
    \item \href{https://en.wikipedia.org/wiki/Humboldt_Prize}{Humboldt Prize} in Germany (2004)
    \item \href{https://en.wikipedia.org/wiki/French_Legion_of_Honor}{French Legion of Honor} in 2010.
    \item \href{https://en.wikipedia.org/wiki/American_Mathematical_Society}{American Mathematical Society} Fellowship (2012).[3]
    \item Foreign honorary member of the \href{https://en.wikipedia.org/wiki/American_Academy_of_Arts_and_Sciences}{American Academy of Arts and Sciences}, 2013.[4]
\end{itemize}

\section{Articles}
\begin{itemize}
    \item Berestycki, Henri; Roquejoffre, Jean-Michel; Rossi, Luca; The influence of a line with fast diffusion on Fisher-KPP propagation. \textit{J. Math. Biol.} 66 (2013), no. 4-5, 743--766.
    \item Barthélemy, Marc; Nadal, Jean-Pierre; Berestycki, Henri Disentangling collective trends from local dynamics. \textit{Proc. Natl. Acad. Sci. USA} 107 (2010), no. 17, 7629--7634.
    \item Berestycki, Henri; Hamel, François; Nadirashvili, Nikolai Elliptic eigenvalue problems with large drift and applications to nonlinear propagation phenomena. \textit{Comm. Math. Phys.} 253 (2005), no. 2, 451--480.
    \item Berestycki, Henri; Hamel, François Front propagation in periodic excitable media. \textit{Comm. Pure Appl. Math.} 55 (2002), no. 8, 949--1032.
    \item Berestycki, H.; Caffarelli, L. A.; Nirenberg, L. Inequalities for second-order elliptic equations with applications to unbounded domains. I. A celebration of John F. Nash, Jr. \textit{Duke Math. J.} 81 (1996), no. 2, 467--494.
    \item Berestycki, H.; Nirenberg, L.; Varadhan, S. R. S. The principal eigenvalue and maximum principle for 2nd-order elliptic operators in general domains. \textit{Comm. Pure Appl. Math.} 47 (1994), no. 1, 47--92.
    \item Berestycki, H.; Lions, P.-L. Nonlinear scalar field equations. I. Existence of a ground state. \textit{Arch. Rational Mech. Anal.} 82 (1983), no. 4, 313--345; II. Existence of infinitely many solutions, \textit{Arch. Rational Mech. Anal.} 82 (1983), no. 4, 347--375.
    \item Bahri, Abbas; Berestycki, Henri A perturbation method in critical point theory and applications. \textit{Trans. Amer. Math. Soc.} 267 (1981), no. 1, 1--32.\hfill$\square$
\end{itemize}

%------------------------------------------------------------------------------%

\chapter{\href{https://en.wikipedia.org/wiki/Ngo_Bao_Chau}{Wikipedia/Ngô Bảo Châu}}

\textbf{Ngô Bảo Châu.}
\begin{itemize}
    \item \textbf{Born.} Jun 28, 1972 (age 48)
    \href{https://en.wikipedia.org/wiki/Hanoi}{Hanoi}, \href{https://en.wikipedia.org/wiki/Democratic_Republic_of_Vietnam}{Democratic Republic of Vietnam}
    \item \textbf{Citizenship.}
    \begin{itemize}
        \item Vietnam
        \item France
    \end{itemize}
    \item \textbf{Alma mater.}
    \begin{itemize}
        \item \href{https://en.wikipedia.org/wiki/%C3%89cole_Normale_Sup%C3%A9rieure}{École Normale Supérieure}
        \item \href{https://en.wikipedia.org/wiki/Universit%C3%A9_de_Paris-Sud}{Université de Paris-Sud}
    \end{itemize}
    \item \textbf{Known for.} Proof of the \href{https://en.wikipedia.org/wiki/Fundamental_lemma_(Langlands_program)}{fundamental lemma}.
    \item \textbf{Spouse(s).} Nguyen Bao Thanh
    \item \textbf{Children.} 3
    \item \textbf{Awards.}
    \begin{itemize}
        \item \href{https://en.wikipedia.org/wiki/Clay_Research_Award}{Clay Research Award} (2004)
        \item \href{https://en.wikipedia.org/wiki/Oberwolfach_Prize}{Oberwolfach Prize} (2007)
        \item \href{https://en.wikipedia.org/wiki/Sophie_Germain_Prize}{Sophie Germain Prize} (2007)
        \item \href{https://en.wikipedia.org/wiki/Fields_Medal}{Fields Metal} (2010)
        \item \href{https://en.wikipedia.org/wiki/Legion_of_Honour}{Legion of Honor} (2011)
    \end{itemize}
\end{itemize}
\textbf{Scientific career.}
\begin{itemize}
    \item \textbf{Fields.} Mathematics
    \item \textbf{Institutions.}
    \begin{itemize}
        \item \href{https://en.wikipedia.org/wiki/Universit%C3%A9_de_Paris-Sud}{Universit\'e de Paris-Sud}
        \item \href{https://en.wikipedia.org/wiki/Institute_for_Advanced_Study}{Institute for Advanced Study}
        \item \href{https://en.wikipedia.org/wiki/University_of_Chicago}{University of Chicago}
    \end{itemize}
    \item \textbf{Doctoral advisor.} \href{https://en.wikipedia.org/wiki/G%C3%A9rard_Laumon}{G\'erard Laumon}
\end{itemize}
\textit{Ngô Bảo Châu} (born Jun 28, 1972)[3] is a Vietnamese-French mathematician at the \href{https://en.wikipedia.org/wiki/University_of_Chicago}{University of Chicago}, best known for proving the \href{https://en.wikipedia.org/wiki/Fundamental_lemma_(Langlands_program)}{fundamental lemma for automorphic forms} proposed by \href{https://en.wikipedia.org/wiki/Robert_Langlands}{Robert Langlands} and \href{https://en.wikipedia.org/wiki/Diana_Shelstad}{Diana Shelstad}.

He is the 1st Vietnamese national to have received the \href{https://en.wikipedia.org/wiki/Fields_Medal}{Fields Medal}.[New Scientist Mathematics `Nobel' rewards boundary-busting work 19 August 2010 ``\textit{Aside from Lindenstrauss, this year's winners were Ngô Bảo Châu of the University of Paris-South, France, Stanslav Smirnov of the University of Geneva, Switzerland, and Cédric Villani of the Henri Poincaré Institute, Paris, France}.''][The Australian Mathematical Society Asia Pacific Mathematics Newsletter April 2011 (pdf) Interview ``\textit{Vietnamese Mathematician Ngô Bảo Châu - From A Mathematical Olympiad Medallist to A Fields Medallist}'' pp. 25--30][Hàm Châu (2005-11-18). ``\textit{Hiện tượng Ngô Bảo Châu}''. Tuổi trẻ Online. Archived from the original on 2010-03-01. Retrieved 2011-10-09.][K.Hưng (2005-12-29). ``\textit{10 sự kiện khoa học — công nghệ nổi bật năm 2005}''. Tuổi trẻ Online. Archived from the original on 2010-02-17. Retrieved 2010-12-19.][Hạ Anh \& Hương Giang, ``\textit{GS Griffiths: `Trong giới Toán học, anh Châu vẫn là người Việt'}'', Vietnamnet. Retrieved 2010-8-19.]

\section{Early life}
Ngô Bảo Châu was born in 1972, the only son of an intellectual family in \href{https://en.wikipedia.org/wiki/Hanoi}{Hanoi}, \href{https://en.wikipedia.org/wiki/North_Vietnam}{North Vietnam}.

His father, professor Ngô Huy Cẩn, is full professor of physics at the Vietnam National Institute of Mechanics.

His mother, Trần Lưu Vân Hiền, is a physician and associate professor at an herbal medicine hospital in Hanoi.

%
The beginning of Châu's schooling was at an experimental elementary school that had been founded by the revolutionary pedagogue Hồ Ngọc Đại, but when his father returned from the Soviet Union with his doctoral degree, he decided that Châu would learn more in traditional schools and enrolled him in the ``chuyên toán'' (special classes for gifted students in mathematics) at the Trưng Vương Middle School.[Koblitz, Neal (2011), ``Interview with Professor Ngô Bảo Châu'', \textit{The Mathematical Intelligencer}, 33 (1): 46--50, doi:10.1007/s00283-010-9184-1, S2CID 120960918]

At age 15, Châu entered the special mathematics class at the \href{https://en.wikipedia.org/wiki/High_School_for_Gifted_Students,_Hanoi_University_of_Science}{High School for Gifted Students, Hanoi University of Science} (Khối chuyên Tổng Hợp -- Đại học Khoa Học Tự Nhiên Hà Nội[``\textit{Đại học Khoa Học Tự Nhiên Hà Nội}''. HUS-VNU official website. Retrieved 2010-08-19.]), formerly known as the A0-class.

In grades 11 and 12, Châu participated in the 29th and 30th \href{https://en.wikipedia.org/wiki/International_Mathematical_Olympiad}{International Mathematical Olympiads} (IMO) and became the 1st Vietnamese student to win 2 IMO gold medals,[11] of which the 1st one was won with a perfect score (42/42).[Ham Chau (2009-02-15). ``Ngô Bao Châu, sommité mondiale des maths''. \textit{Le Courrier du Vietnam} (in French).]

%
After high school, Châu expected to study in \href{https://en.wikipedia.org/wiki/Budapest}{Budapest}, but in the aftermath of the fall of \href{https://en.wikipedia.org/wiki/Communism}{Communism} in \href{https://en.wikipedia.org/wiki/Eastern_Europe}{Eastern Europe}, the new \href{https://en.wikipedia.org/wiki/Hungary}{Hungarian} government halted scholarships to students from Vietnam.[Hàm Châu (2010-08-17). ``Ngô Bảo Châu, ``bom tấn'' và ``trống đồng'' trong toán học'' (in Vietnamese). Dân Trí.]

After visiting Châu's father, Paul Germain, secretary of the \href{https://en.wikipedia.org/wiki/French_Academy_of_Sciences}{French Academy of Sciences}, arranged for Châu to study in France.

He was offered a scholarship by the French government for undergraduate study at the \href{https://en.wikipedia.org/wiki/Paris_VI_University}{Paris VI University}, then in 1992, he entered the \href{https://en.wikipedia.org/wiki/%C3%89cole_Normale_Sup%C3%A9rieure}{École Normale Supérieure}.

He obtained a \href{https://en.wikipedia.org/wiki/Doctor_of_Philosophy}{PhD} in 1997 from the \href{https://en.wikipedia.org/wiki/Universite_Paris-Sud}{Universite Paris-Sud} under the supervision of \href{https://en.wikipedia.org/wiki/G%C3%A9rard_Laumon}{Gérard Laumon}.
    
He became a member of \href{https://en.wikipedia.org/wiki/CNRS}{CNRS} at \href{https://en.wikipedia.org/wiki/Paris_13_University}{Paris 13 University} from 1998 to 2005, and defended his \href{https://en.wikipedia.org/wiki/Habilitation}{habilitation} degree there in 2003.

He holds both Vietnamese and French citizenship.[14]

\section{Career}
Châu became a professor at \href{https://en.wikipedia.org/wiki/Paris-Sud_11_University}{Paris-Sud 11 University} in 2005.

In 2005, at age 33, Chau received the title of professor in Vietnam, becoming the country's youngest-ever professor.[Ham Chau (2009-02-15). ``Ngô Bao Châu, sommité mondiale des maths''. \textit{Le Courrier du Vietnam} (in French).]

Since 2007, Châu has worked at the \href{https://en.wikipedia.org/wiki/Institute_for_Advanced_Study}{Institute for Advanced Study}, \href{https://en.wikipedia.org/wiki/Princeton,_New_Jersey}{Princeton, New Jersey}, as well as the Hanoi Institute of Mathematics.[``\textit{Top 10 Scientific Discoveries of 2009}''. Time magazine. 2009-12-08.]

He joined the mathematics faculty at the \href{https://en.wikipedia.org/wiki/University_of_Chicago}{University of Chicago} on Sep 1, 2010.

In addition, since 2011 he has been Scientific Director of the newly founded Vietnam Institute for Advanced Study in Mathematics (VIASM).[``History''. Vietnam Institute for Advanced Study in Mathematics. Retrieved Mar 15, 2019.]

In 2016 Châu was Co-General Chair of \href{https://en.wikipedia.org/wiki/Asiacrypt}{Asiacrypt} the 1st time that the Asian cryptography conference was held in Vietnam.[17]

\section{Work}
Châu 1st came to prominence by proving, in joint work with \href{https://en.wikipedia.org/wiki/G%C3%A9rard_Laumon}{Gérard Laumon}, the \href{https://en.wikipedia.org/wiki/Fundamental_lemma_(Langlands_program)}{fundamental lemma} for \href{https://en.wikipedia.org/wiki/Unitary_group}{unitary groups}.
    
Their general strategy was to understand the local \href{https://en.wikipedia.org/wiki/Orbital_integral}{orbital integrals} appearing in the fundamental lemma in terms of \textit{affine Springer fibers} arising in the \href{https://en.wikipedia.org/wiki/Hitchin_fibration}{Hitchin fibration}.

This allowed them to employ the tools of \textit{geometric representation theory}, namely the theory of \href{https://en.wikipedia.org/wiki/Perverse_sheaves}{perverse sheaves}, to study what was initially a \href{https://en.wikipedia.org/wiki/Combinatorics}{combinatorial} problem of a \href{https://en.wikipedia.org/wiki/Number_theory}{number-theoretic} nature.

Chau eventually succeeded in formulating the proof for the fundamental lemma for \href{https://en.wikipedia.org/wiki/Lie_algebra}{Lie algebras} in 2008.[12]

Together with results from \href{https://en.wikipedia.org/wiki/Jean-Loup_Waldspurger}{Jean-Loup Waldspurger}, who had earlier deduced stronger forms of the fundamental lemma from this result, this completed the proof of the fundamental lemma in all cases.

As a result, Châu was awarded a \href{https://en.wikipedia.org/wiki/Fields_Medal}{Fields Medal} in 2010.

%
Ngo Bao Chau was the co-author of the Vietnamese children book ``Ai and Ky in the land of the invisible numbers''.

\section{Honors}
In 2004, Chau and Laumon were awarded the \href{https://en.wikipedia.org/wiki/Clay_Research_Award}{Clay Research Award} for their achievement in solving the fundamental lemma proposed by \href{https://en.wikipedia.org/wiki/Robert_Langlands}{Robert Langlands} for the case of \href{https://en.wikipedia.org/wiki/Unitary_group}{unitary groups}.[12]

Chau's proof of the general case was selected by \href{https://en.wikipedia.org/wiki/Time_(magazine)}{\textit{Time}} as 1 of the Top Ten Scientific Discoveries of 2009.[``\textit{Top 10 Scientific Discoveries of 2009}''. Time magazine. 2009-12-08.]

In 2010, he received the \href{https://en.wikipedia.org/wiki/Fields_Medal}{Fields Medal} and in 2011, he received the \href{https://en.wikipedia.org/wiki/Legion_of_Honour}{Legion of Honour}.[18]

In 2012 he became a fellow of the \href{https://en.wikipedia.org/wiki/American_Mathematical_Society}{American Mathematical Society}.[19]

\section{Major publications}
\begin{itemize}
    \item Ngô, Bao Châu. \textit{Le lemme fondamental pour les algèbres de Lie}. [\textit{The fundamental lemma for Lie algebras}] Publ. Math. Inst. Hautes Études Sci. 111 (2010), 1--169. doi:10.1007/s10240-010-0026-7 MR2653248 arXiv:0801.0446
\end{itemize}

\section{External links}
\begin{itemize}
    \item \href{http://www.math.uchicago.edu/~ngo/nbc-homepage.html}{Ngô Bảo Châu's homepage at the University of Chicago}
    \item \href{https://web.archive.org/web/20130102171548/http://ngobaochau.wordpress.com/}{Personal blog} (in Vietnamese)
    \item \href{https://arxiv.org/abs/math/0404454}{Ngo Bao Chau and Gérard Laumon's proof of the Fundamental Lemma for unitary groups.}
    \item \href{https://web.archive.org/web/20090306143328/http://www.mfo.de/programme/prize/Ngo2008.pdf}{The Oberwolfach Prize 2007 is awarded to Ngo Bao Chau (Laudatio by Prof. Rapoport)}\hfill$\square$
\end{itemize}

%------------------------------------------------------------------------------%

\part{Physicists}

%------------------------------------------------------------------------------%

\part{Computer Scientists}

%------------------------------------------------------------------------------%

\part{Psychologists}

%------------------------------------------------------------------------------%

\part{Philosophers}

%------------------------------------------------------------------------------%

\printbibliography[heading=bibintoc]
\end{document}